\section{Double counting and the handshake theorem}\label{ch07:sec:double}
Add up the degrees in the example graph $G$, which has four edges: $2+2+3+1=8=2\cdot4$. This is no accident.

Consider every pair $(v,e)$ in which $v$ is a vertex, $e$ is an edge, and $v$ is an endpoint of $e$, that is, $v\in e$. The table below has one row for each vertex of $G$ and one column for each edge; an entry $1$ marks a pair $(v,e)$ with $v\in e$.
\begin{center}
\begin{tabular}{@{}lccccc@{}}
\toprule
 & $ab$ & $bc$ & $ca$ & $cd$ & row total\\
\midrule
$a$ & $1$ & $0$ & $1$ & $0$ & $2$\\
$b$ & $1$ & $1$ & $0$ & $0$ & $2$\\
$c$ & $0$ & $1$ & $1$ & $1$ & $3$\\
$d$ & $0$ & $0$ & $0$ & $1$ & $1$\\
\midrule
column total & $2$ & $2$ & $2$ & $2$ & $8$\\
\bottomrule
\end{tabular}
\end{center}
Count the entries equal to $1$ in two ways. The row of a vertex has a $1$ for each edge at that vertex, so the row totals are the degrees, which add up to $8$. Every edge has two endpoints, so every column total is $2$, and the columns also add up to $8$: both totals count the same entries.

Counting one collection of objects in two ways and setting the two answers equal is called \emph{double counting}. Its essential step is to describe the collection precisely, so that we can check that each grouping counts every member exactly once.

To state the general result, we extend the sum notation of Section~\ref{ch01:sec:notation}: $\sum_{v\in V}\deg(v)$ means that we add $\deg(v)$ once for each vertex $v$ in $V$; for the example graph it is $\deg(a)+\deg(b)+\deg(c)+\deg(d)=8$. The order of the terms does not matter, because the rules $a+b=b+a$ and $(a+b)+c=a+(b+c)$ of Section~\ref{ch01:sec:letters} let us reorder and regroup the terms of any finite sum. In the same way, $\sum_{e\in E}2$ adds $2$ once for each edge, so it equals $2|E|$.

\par\noindent\begin{minipage}{\linewidth}
\begin{theorem}[Handshake theorem]\label{thm:handshake}
In every graph with vertex set $V$ and edge set $E$,
\[
\sum_{v\in V}\deg(v)=2|E|.
\]
\end{theorem}
\end{minipage}\par

The name comes from a party: if the guests are vertices and the handshakes are edges, then $\deg(v)$ is the number of hands that $v$ shook, and each handshake is counted twice in the total, once by each of the two people involved.

\begin{proof}
Let $P$ be the set of all pairs $(v,e)$ with $v\in V$, $e\in E$, and $v\in e$. We count $P$ in two ways.

First group the pairs by their vertex. For a fixed vertex $v$, the pairs in $P$ with first entry $v$ are the pairs $(v,e)$ in which $e$ is an edge at $v$, and by Section~\ref{ch07:sec:graph} there are exactly $\deg(v)$ such edges. Each pair in $P$ has just one first entry, so it lies in exactly one of these groups, and $|P|=\sum_{v\in V}\deg(v)$.

Next group the pairs by their edge. For a fixed edge $e=\{u,w\}$, the pairs in $P$ with second entry $e$ are $(u,e)$ and $(w,e)$, two different pairs because $u\ne w$. Again every pair lies in exactly one group, so $|P|=\sum_{e\in E}2=2|E|$.

Both expressions equal $|P|$, so they are equal to each other.
\end{proof}

In the example graph, two vertices, $c$ and $d$, have odd degree. That this number is even is no coincidence either.

\par\noindent\begin{minipage}{\linewidth}
\begin{proposition}\label{ch07:prop:odd}
In every graph, the number of vertices of odd degree is even.
\end{proposition}
\end{minipage}\par

\begin{proof}
Let $k$ be the number of vertices of odd degree. By Lemma~\ref{ch04:lem:even-or-odd}, every degree is either even or odd, and not both. So each vertex contributes its degree to exactly one of two sums: $S_{\mathrm{even}}$, over the vertices of even degree, and $S_{\mathrm{odd}}$, over the vertices of odd degree. Together they make up the degree sum, so $S_{\mathrm{even}}+S_{\mathrm{odd}}=2|E|$ by the handshake theorem.
\begin{itemize}
\item $S_{\mathrm{even}}$ is even, because a factor $2$ can be taken out of every term. (If no vertex has even degree, $S_{\mathrm{even}}$ is the empty sum $0$, which is even too.)
\item Therefore $S_{\mathrm{odd}}=2|E|-S_{\mathrm{even}}$ is even, since $2x-2y=2(x-y)$.
\end{itemize}
Now write the degrees of the $k$ vertices of odd degree, using Definition~\ref{def:parity}, as $2\ell_1+1,\ldots,2\ell_k+1$ with integers $\ell_1,\ldots,\ell_k$. Adding them gives
\[
S_{\mathrm{odd}}=(2\ell_1+1)+\cdots+(2\ell_k+1)=2(\ell_1+\cdots+\ell_k)+k,
\]
because each of the $k$ terms contributes one $1$. Hence $k=S_{\mathrm{odd}}-2(\ell_1+\cdots+\ell_k)$ is the difference of two even numbers, so $k$ is even. (If $k=0$, the conclusion holds at once, since $0$ is even.)
\end{proof}

For instance, the proposition shows that seven people at a party cannot each shake hands with exactly three of the others: all seven vertices would have the odd degree $3$, and $7$ is odd. (The handshake theorem shows this directly: the degree sum $7\cdot3=21$ would be odd, but it equals $2|E|$.)
